\section{1. Scope and scientific
question}\label{scope-and-scientific-question}

The composition- and temperature-dependent zone-center gap of
\(\mathrm{Hg}_{1-x}\mathrm{Cd}_x\mathrm{Te}\) is central to infrared
cutoff wavelength, carrier statistics, transport, and optical-response
modeling. Empirical gap relations are often compared at sub-meV to
few-meV resolution \citep{Hansen1982, Seiler1990, Laurenti1990}. Such
comparisons are meaningful only if the experimentally reported fitted or
operational coordinate is defined with comparable rigor.

An optical absorption curve contains no unique model-independent point
labeled \(E_g\). A reported edge can be generated by a power-law
intercept, a Kane-region fit, an Urbach-tail intersection, a derivative
construction, a fixed absorption crossing, or a detector-response
criterion. These procedures can be useful, but they are not
automatically estimators of the same measurand.

This paper addresses a narrow methodological question:

\begin{quote}
For two published HgCdTe absorption-coefficient curves, how strongly
does the fitted optical intercept depend on the declared extraction
model, reconstruction rule, fit-domain endpoints, and point weighting?
\end{quote}

The paper does \textbf{not} attempt to determine the true band gap of
either specimen or rank universal gap equations. It is a two-spectrum
case study of downstream extraction dependence.

\begin{figure}
\centering
\includegraphics{figures/figure1_inference_chain.pdf}
\caption{Inference chain from specimen state to a reported optical
coordinate.}
\end{figure}
